
% 2. 基本グラフィック・色・レイアウト
\usepackage{fancybox}
\usepackage{graphicx,color}
\usepackage[margin=20mm]{geometry}
\usepackage{multicol,multirow,varwidth,float,caption}

% 3. emathコアパッケージ
\usepackage[deepenum,notMy]{emath}

% emath拡張パッケージ群
\usepackage{hako}        % 穴埋め（箱）
\usepackage{EMfbox}      % 囲み枠
\usepackage{abcgreek}    % ギリシャ文字
\usepackage{emathEy}     % 記号・特殊文字
\usepackage{emathG}      % 幾何（図形）
\usepackage{emathT}      % 表作成
\usepackage{emathPs}     % PostScript連携
\usepackage{emathR}      % 実数計算
\usepackage{emathAe}     % 解答作成
\usepackage{emathThm}    % 定理環境
\usepackage{EMpict2e}    % 描画拡張
\usepackage[deluxe]{emathOtf} % OTFフォント（太字などの対応）


% 5. 各種定義と書式設定
% 長丸囲み文字の定義
\def\nagamaru#1{\mbox{\pIIebNagamaru{#1}}}

% 定理環境（「仮定」など）の共通テスト風デザイン
\theorembodyfont{\normalfont}
\theoremstyle{boxed}
\newtheorem<frame=rectbox,frameoption={oval=8pt,preitem=~,postitem=~}>{katei}{仮定}

% キャプションの日本語化（図1: → 図 1）
\captionsetup{figurename=図,labelsep=quad}

% カウンターのリセット設定（大問ごとに数式や図の番号をリセット）
\resetcounter{equation}[enumi]
\resetcounter{figure}[enumi]
\resetcounter{table}[enumi]
\EMfigtrue
\apnlist{\setlength{\listparindent}{1zw}}

% 大問ラベルの書式定義（専用の箇条書き環境として定義）
\newenvironment{QuestionEnum}{%
    \begin{enumerate}<apnlist={\leftmargin1zw\itemindent4zw}>'{\bfseries\Large 第\arabic{enumi}問}'%
}{%
    \end{enumerate}%
}


% --- 単線枠 ---
\newcommand{\fgb}[1]{%
    {% 
    \setlength{\fboxrule}{1.6pt}%
    \framebox[1.7cm]{#1}%
    }%
}
\newcommand{\fb}[1]{%
    {% 
    \setlength{\fboxrule}{0.4pt}%
    \framebox[1.7cm]{#1}%
    }%
}

% --- 二重線枠 ---
\newcommand{\dgb}[1]{%
    {%
    \edef\originalsep{\the\fboxsep}% 現在の標準の余白サイズを記憶
    \setlength{\fboxrule}{1.6pt}% ① 外側の線を太くする(1.6pt)
    \setlength{\fboxsep}{1.6pt}%  ② 二重線の「線と線のすき間」の幅
    \fbox{%
        \setlength{\fboxrule}{0.4pt}% ③ 内側の線は普通の太さ(0.8pt)
        \fbox{\makebox[1.4cm]{#1}}%
    }%
    }%
}
\newcommand{\db}[1]{%
    {%
    \edef\originalsep{\the\fboxsep}% 現在の標準の余白サイズを記憶
    \setlength{\fboxrule}{0.4pt}% ① 外側の線を太くする(1.6pt)
    \setlength{\fboxsep}{1.6pt}%  ② 二重線の「線と線のすき間」の幅
    \fbox{%
        \setlength{\fboxrule}{0.4pt}% ③ 内側の線は普通の太さ(0.8pt)
        \fbox{\makebox[1.4cm]{#1}}%
    }%
    }%
}

\newcommand{\subject}[3]{%
\begin{center}
    {\LARGE \textbf{#1}}\\[0.5em]
    {\Huge （}%
    \raisebox{0.4em}{%
        {\normalsize \textbf{解答番号}} \fb{#2} ～ \fb{#3}%
    }%
    {\Huge ）}\\
\end{center}
}